\chapter{Ví dụ minh hoạ trình bày nội dung}
\label{ch:noi_dung_mau}

% Đây là chương MẪU minh hoạ cách chèn hình, bảng, code, công thức, trích dẫn.
% Khi viết chương thật (chương 2, 3, 4... tuỳ đề tài) hãy copy các ví dụ này
% và xoá phần không cần.

\section{Chèn hình ảnh}
\begin{figure}[H]
    \centering
    % \includegraphics[width=0.7\linewidth]{images/vi_du.png}
    \todo{Chèn ảnh vào đây, ví dụ: \texttt{images/vi\_du.png}}
    \caption{Ví dụ chú thích hình}
    \label{fig:vi_du}
\end{figure}

Tham chiếu tới hình bằng lệnh \texttt{\textbackslash ref\{fig:vi\_du\}}, ví dụ: xem Hình~\ref{fig:vi_du}.

\section{Chèn bảng}
\begin{table}[H]
\centering
\caption{Ví dụ bảng số liệu}
\label{tab:vi_du}
\begin{tabular}{lll}
\toprule
\textbf{Cột 1} & \textbf{Cột 2} & \textbf{Cột 3} \\
\midrule
Dữ liệu A & Dữ liệu B & Dữ liệu C \\
Dữ liệu D & Dữ liệu E & Dữ liệu F \\
\bottomrule
\end{tabular}
\end{table}

\section{Chèn đoạn code}
\begin{lstlisting}[language=Python, caption={Ví dụ đoạn code minh hoạ}]
def hello():
    print("Xin chao bao cao!")
\end{lstlisting}

\section{Công thức toán}
Ví dụ một công thức đơn giản:
\begin{equation}
    E = mc^2
    \label{eq:vi_du}
\end{equation}

\section{Trích dẫn tài liệu tham khảo}
Đây là ví dụ trích dẫn tài liệu tham khảo \cite{example2024}.
